% Generated draft: requires mathematical and source review.
\documentclass[11pt]{article}
\usepackage[margin=25mm]{geometry}
\usepackage{fontspec}
\setmainfont{Latin Modern Roman}
\usepackage{amsmath,amssymb,mathtools,unicode-math}
\setmathfont{Latin Modern Math}
\usepackage{enumitem,needspace,xurl,hyperref}
\setlength{\parindent}{0pt}
\setlength{\parskip}{5pt}
\setlength{\emergencystretch}{3em}
\newcommand{\N}{\mathbb{N}}
\newcommand{\Z}{\mathbb{Z}}
\newcommand{\Q}{\mathbb{Q}}
\newcommand{\R}{\mathbb{R}}
\newcommand{\C}{\mathbb{C}}
\newcommand{\F}{\mathbb{F}}
\newcommand{\A}{\mathbb{A}}
\newcommand{\PP}{\mathbb{P}}
\newcommand{\E}{\mathbb{E}}
\newcommand{\eps}{\varepsilon}
\newcommand{\dd}{\,\mathrm{d}}
\begin{document}
\section*{UnsolvedMath ID: 20002453}
\textbf{Source problem code:} catalogue entry
\subsection{Definitions and mathematical statement}
Let $N\ge3$ and consider $u:\R^N\times\R\to\R^M$ solving either
\[
 J u_t=\Delta u+g(u)\quad\text{(NLS)},\qquad
 u_{tt}=\Delta u+g(u)\quad\text{(NLKG)},
\]
where $J$ is invertible and skew-symmetric and the source assumes continuous radial $g(y)=h(|y|^2)y/|y|$ for $y\ne0$, continuously extended at $0$. Warchall's standing waves are
\[
 u(x,t)=e^{\mu Kt}\Phi(x),\qquad
 \Phi(x)=\widehat\psi(x/|x|)w(|x|),
\]
where $K$ is real skew-symmetric, $K=J^{-1}$ for NLS, $K^2=-I$ for NLKG, and
\[
 \widehat\psi:S^{N-1}\to S^{M-1},\qquad
 \Delta_S\widehat\psi=-l(l+N-2)\widehat\psi.
\]
The possible angular degrees are restricted by the target dimension $M$. The question asks under what conditions these encapsulated vortices are stable or unstable. The source cites localized profiles under suitable hypotheses including $f'(0)<0$ and $\int_0^a f(s)\,\dd s>0$ for some $a>0$, and notes Galilean or Lorentz boosts as routes to traveling waves; it does not provide a complete existence or stability hypothesis list.

There is a sign conflict in the supplied setup. For the literal NLS equation, $K=J^{-1}$ and $e^{+\mu Kt}$ imply
\[
 \Delta\Phi+g(\Phi)+\omega\Phi=0,
 \qquad\omega=-\mu,
\]
whereas the source prints $\omega=+\mu$. For NLKG the same substitution consistently gives $\omega=\mu^2$. Thus the radial profile equation is
\[
 w''+\frac{N-1}{r}w'-\frac{l(l+N-2)}{r^2}w+f(w)=0,
\]
with $f(w)=\operatorname{sgn}(w)h(w^2)+\omega w$ for $w\ne0$, interpreted through its continuous extension. The source's NLS phase signs cannot all be imposed simultaneously.

The aid proves only a conditional linear high-angular-sector filter. Assume now $g$ is $C^1$ and a radial gradient, the chosen profile is smooth and localized, and
\[
 L=-\Delta-Dg(\Phi)-\omega I,
 \quad -Dg(0)-\omega I\ge\gamma I\ (\gamma>0),\quad
 C=\operatorname*{ess\,sup}_{x\ne0}
 |x|^2\max\{\lambda_{\max}(Dg(\Phi(x))-Dg(0)),0\}<\infty.
\]
Assume also a compact group of combined spatial and orthogonal target symmetries fixes $\Phi$, commutes with the evolution matrices, and preserves the operators. In one invariant isotypic sector $\sigma$, let $\Lambda_\sigma$ be an angular lower bound, meaning every angular section in that sector satisfies
\[
 \int_{S^{N-1}}|\nabla_Sv|^2
 \ge\Lambda_\sigma\int_{S^{N-1}}|v|^2.
\]
Then $\Lambda_\sigma>C$ yields the reported positive quadratic-form bound
\[
 \langle Lv,v\rangle\ge
 \|\partial_rv\|_2^2
 +(\Lambda_\sigma-C)\|r^{-1}v\|_2^2
 +\gamma\|v\|_2^2.
\]
For the co-rotating linearizations
\[
 Jv_t=-Lv\quad\text{(NLS)},\qquad
 v_{tt}+2\mu Kv_t=-Lv\quad\text{(NLKG)},
\]
this gives conserved positive quadratic energies, respectively $\tfrac12\langle Lv,v\rangle$ and $\tfrac12(\|v_t\|_2^2+\langle Lv,v\rangle)$, excluding finite-energy exponential growth within that sector. If only finitely many sectors have $\Lambda_\sigma\le C$, all possible exponential instabilities are confined to them. This finiteness condition is an additional hypothesis. Scalar spherical-harmonic degrees need not separately be invariant, since the angular dependence of $Dg(\Phi)$ can couple them; combined isotypic sectors retain their multiplicities. The aid does not prove stability of the remaining low sectors or nonlinear orbital stability of the full vortex.
\subsection{Short English statement}
Determine stability of higher-dimensional encapsulated vortices in vector NLS or NLKG. The supplied aid identifies a source sign conflict and, under additional smoothness, gap and symmetry hypotheses, excludes linear exponential growth in sufficiently high combined angular sectors; low sectors and nonlinear stability remain to be analyzed.
\subsection{Sources}
\begin{itemize}
\item H. Warchall, ``Stability of Encapsulated-Vortex Solutions,'' American Institute of Mathematics workshop \emph{Stability criteria for multi-dimensional waves and patterns}, page 7 as identified by supplied verification, \url{https://aimath.org/WWN/multidimwaves/multidimwaves.pdf}.
\item Local UnsolvedMath registry, permanent ID 20002453, source code \texttt{catalogue entry}, canonical location \texttt{aim-pdes-notes.json}, \texttt{notes[64]}, record dated 14 August 2026. Companion permanent ID 20002452, \texttt{AIM-PDES-0064}, \texttt{notes[63]}, contains the setup lost across the record boundary and the following supplied references.
\item J. Iaia and H. A. Warchall, ``Encapsulated-vortex solutions to equivariant wave equations: existence,'' \emph{SIAM J. Math. Anal.} 30 (1999), 118--139.
\item R. L. Pego and H. A. Warchall, ``Spectrally stable encapsulated vortices for nonlinear Schr\"odinger equations,'' \emph{J. Nonlinear Sci.} 12 (2002), 347--394.
\item Supplied research attempt 1: partial result, problem status at run partially solved, rigor labelled proved, novelty confidence low. Its sectorwise coercivity filter is machine-generated research aid, not peer reviewed, requiring independent verification before citation; the original literature field is empty despite the references embedded in the companion statement.
\end{itemize}
\end{document}
